% Part: first-order-logic
% Chapter: model-theory
% Section: theory-of-m

\documentclass[../../../include/open-logic-section]{subfiles}

\begin{document}

\olfileid{mod}{bas}{thm}

\section{Teori Sawijining \printtoken{S}{structure}}

Saben !!{structure}~$\Struct{M}$ ndadekake sawenehing
!!{sentence}s bener lan sawenehing liyane kleru.
Himpunan kabeh !!{sentence}s kang digawe bener
dening struktur iku diarani \emph{teori}-ne.
Himpunan iku pancen sawijining teori, amarga
apa wae kang dadi akibat semantise mesthi
bener ing kabeh modele, kalebu~$\Struct{M}$.

\begin{defn}
  Yen diwenehi !!a{structure}~$\Struct M$,
  \emph{teori} saka $\Struct{M}$ yaiku
  himpunan $\Theory{M}$ kang isi !!{sentence}s
  kang bener ing $\Struct{M}$, yaiku
  $\Theory{M} = \Setabs{!A}{\Sat{M}{!A}}$.
\end{defn}

Kita uga nggunakake tembung teori
kanthi ora resmi kanggo nyebut himpunan
!!{sentence}s kang duwe interpretasi
kang dikarepake, apa wis katutup
miturut deduksi utawa durung.

\begin{prop}
Kanggo sembarang $\Struct{M}$,
$\Theory{M}$ iku komplet.
\end{prop}

\begin{proof}
Kanggo saben !!{sentence}~$!A$,
bisa $\Sat{M}{!A}$ utawa
$\Sat{M}{\lnot !A}$,
dadi bisa $!A \in \Theory{M}$
utawa $\lnot !A \in \Theory{M}$.
\end{proof}

\begin{prop}\ollabel{prop:equiv}
  Yen $\Struct{N} \models !A$
  kanggo saben $!A \in \Theory{M}$,
  mula $\Struct{M} \elemequiv \Struct{N}$.
\end{prop}

\begin{proof}
Amarga $\Sat{N}{!A}$ kanggo kabeh
$!A \in \Theory{M}$, mula
$\Theory{M} \subseteq
\Theory{N}$. Yen $\Sat{N}{!A}$,
mula $\Sat/{N}{\lnot !A}$,
saéngga $\lnot !A
\notin \Theory{M}$. Amarga
$\Theory{M}$ komplet, mula
$!A \in \Theory{M}$. Dadi,
$\Theory{N} \subseteq \Theory{M}$,
lan kita nduweni
$\Struct{M} \elemequiv \Struct{N}$.
\end{proof}

\begin{rem}\ollabel{remark:R}
  Gatekna $\Struct{R} = \langle\Real, <\rangle$,
  yaiku !!{structure} kang ranahé
  himpunan $\Real$ saka wilangan real,
  ing !!{language} kang mung ngemot
  !!{predicate} kanthi rong panggonan
  kang diinterpretasi minangka relasi
  $<$ ing wilangan real. Cetha yen
  $\Struct{R}$ iku !!{nonenumerable};
  nanging amarga $\Theory{R}$
  mesthi konsisten, miturut teorema
  L\"owenheim--Skolem teori iku
  duwe model kang !!a{enumerable},
  upamane $\Struct{S}$, lan
  miturut \olref{prop:equiv},
  $\Struct{R}
  \equiv \Struct{S}$. Sabanjure,
  amarga $\Struct{R}$ lan $\Struct{S}$
  ora isomorfis, iki nuduhake yen
  kosok baline
  \olref[iso]{thm:isom}
  ora mesthi bener.
\end{rem}

\end{document}
